\section{Ciclos históricos de infraestructura: escasez, sobreconstrucción y estandarización}

El ponente sitúa el compute actual dentro de una historia más larga de la infraestructura: el acero, el ferrocarril, la electricidad, la fibra óptica, los materiales semiconductores y las tecnologías energéticas pasaron por un auge de la demanda, una afluencia de capital, una expansión de la oferta, una caída de precios, la consolidación de empresas y la formación de estándares. La analogía histórica sirve para generar preguntas, como «dónde habrá sobreconstrucción», «qué interfaces se estandarizarán» o «quién asume los costos hundidos»; no permite predecir calendarios directamente, porque cada recurso difiere en sus propiedades físicas, su estructura regulatoria y su ritmo de innovación.

\subsection{Gilded Age: la infraestructura y el poder se expanden juntos}

La Gilded Age de finales del siglo XIX suele usarse para describir una etapa de rápida expansión del ferrocarril, el acero y las finanzas, acompañada de concentración del poder de mercado, conflictos laborales y una regulación pública cada vez más intensa. El ponente la utiliza para recordar a los estudiantes que el auge de una tecnología de propósito general no solo trae más capacidad productiva, sino que también redistribuye el poder de negociación. Por ello, la cuestión central de la infraestructura de IA no es solo «cuánto puede construirse», sino también quién controla el acceso, quién fija los precios y quién puede seguir innovando en tiempos de escasez.

\begin{figure}[H]
\centering
\includegraphics[width=0.88\textwidth]{images/00-51-00-gilded-age.jpg}
\caption{La Gilded Age sirve como punto de entrada histórico de esta clase: el auge de la infraestructura y la concentración del poder suelen aparecer al mismo tiempo.}
\figsource{Video oficial de Stanford Online 00:51:00, `images/00-51-00-gilded-age.jpg`.}
\end{figure}

\begin{knowledgebox}{Lectura de la figura: en una analogía histórica, primero se comparan los mecanismos}
Los mecanismos comparables incluyen los altos costos fijos, los efectos de red, el control de los recursos, la integración vertical, la discriminación de precios y la presión de la regulación pública. Las partes que no admiten una analogía directa incluyen la velocidad de actualización tecnológica, las cadenas de suministro globales, el costo de replicar software y los mercados de capital modernos. Tomar «hoy es una nueva Gilded Age» como conclusión es demasiado burdo; tomarlo como un recordatorio para examinar la concentración del mercado y el interés público tiene más valor didáctico.
\end{knowledgebox}

\subsection{El acero: cómo las mejoras de proceso cambian el precio de un recurso}

Innovaciones de proceso como el Bessemer process redujeron de forma notable el costo de producción del acero, lo que permitió que el ferrocarril, la construcción y la maquinaria se expandieran a mayor escala. Este caso corresponde, en la IA, a la eficiencia de los modelos, el diseño de chips y la optimización de sistemas: el crecimiento de la demanda no tiene que satisfacerse únicamente aumentando la oferta de materias primas, porque el propio proceso de producción también modifica el costo unitario. A su vez, la caída de precios libera nueva demanda, de modo que el consumo total sigue aumentando.

\begin{figure}[H]
\centering
\includegraphics[width=0.94\textwidth]{images/00-52-00-bessemer-steel-prices.jpg}
\caption{Tras el proceso Bessemer, el precio del acero bajó durante un largo periodo: la eficiencia productiva cambió la accesibilidad de la infraestructura.}
\figsource{Video oficial de Stanford Online 00:52:00, `images/00-52-00-bessemer-steel-prices.jpg`.}
\end{figure}

\begin{importantbox}{Lectura de la figura: la caída del costo unitario puede ampliar la demanda total}
Cuando el acero se abarató, la sociedad no dejó de usarlo, sino que lo llevó a más vías férreas, puentes y edificios. La eficiencia de cómputo también puede producir un rebound effect similar: si cada inferencia es más barata, surgen más agent, tareas más largas e interacciones más frecuentes. Al estimar la demanda total de energía y de capacidad de cómputo, no basta con ver la reducción del costo por tarea; también hay que estimar la elasticidad de la demanda que generan los nuevos usos.
\end{importantbox}

\subsection{La fibra óptica: la sobreconstrucción también puede dejar capacidad pública}

La gran cantidad de fibra óptica tendida durante la burbuja de internet provocó a corto plazo quiebras y depreciación de activos, pero dejó capacidad de bajo costo para la banda ancha y los servicios en la nube posteriores. Esta historia suele usarse para justificar la expansión de los centros de datos de IA, pero entre ambos casos hay diferencias importantes: la fibra oscura puede permanecer inactiva durante mucho tiempo, mientras que las GPU se deprecian más rápido, y la renovación generacional de los chips, el encapsulado y la compatibilidad del software acortan su vida útil efectiva. Que la sobreconstrucción tenga valor público a largo plazo depende de que los activos puedan reutilizarse.

\begin{figure}[H]
\centering
\includegraphics[width=0.94\textwidth]{images/00-53-00-fiber-optic-overbuild.jpg}
\caption{El precio de la fibra óptica cayó rápidamente tras la sobreconstrucción, y la capacidad ociosa se convirtió después en la base de la expansión de internet.}
\figsource{Video oficial de Stanford Online 00:53:00, `images/00-53-00-fiber-optic-overbuild.jpg`.}
\end{figure}

\begin{warningbox}{Lectura de la figura: ``primero sobreconstruir y luego ver'' no es una estrategia incondicional}
La vida útil de los activos, los costos de mantenimiento, la compatibilidad con los estándares y la dificultad de redistribuirlos determinan si la sobreconstrucción puede convertirse en capacidad pública. Si un centro de datos está atado a una sola red eléctrica, a una interconexión cerrada o a chips que se vuelven obsoletos con rapidez, su valor residual puede ser muy bajo. La planeación de infraestructura debe considerar al mismo tiempo la modularity, la posibilidad de actualización y las interfaces abiertas, en lugar de comparar solo la escala máxima de construcción.
\end{warningbox}

\subsection{La DRAM: por qué la oferta de semiconductores es cíclica}

\term{DRAM} significa Dynamic Random-Access Memory; en español, memoria dinámica de acceso aleatorio. Almacena cada bit en un capacitor, requiere refrescarse periódicamente y se usa ampliamente como memoria principal en servidores y computadoras personales. La inversión en plantas de fabricación de obleas tiene ciclos largos y costos fijos altos, mientras que la demanda varía con la coyuntura de los dispositivos y los centros de datos, por lo que los precios suelen oscilar entre la escasez y el exceso. La memoria de video y el encapsulado avanzado que usan los aceleradores de IA están sujetos a restricciones similares de planeación de capacidad.

\begin{figure}[H]
\centering
\includegraphics[width=0.94\textwidth]{images/00-53-30-dram-cycles.jpg}
\caption{El ciclo de la industria de la DRAM muestra cómo se repiten el retraso en la ampliación de capacidad, el ajuste de inventarios y la volatilidad de precios.}
\figsource{Video oficial de Stanford Online 00:53:30, `images/00-53-30-dram-cycles.jpg`.}
\end{figure}

\begin{knowledgebox}{Lectura de la figura: cómo el retraso de la oferta genera ciclos}
Cuando la demanda aumenta, primero suben los precios y los fabricantes amplían su capacidad en consecuencia; la nueva capacidad entra de golpe tras la construcción y el arranque gradual, el mercado puede pasar al exceso de oferta, y la caída de precios obliga a las empresas a recortar la inversión. Cuando la demanda se recupera, los recortes de la ronda anterior vuelven a estrechar la oferta. Las GPU, la HBM y el encapsulado avanzado no son idénticos a la DRAM, pero en todos los casos hay que analizar en un mismo modelo los ciclos largos de construcción y los pronósticos de demanda de corto plazo.
\end{knowledgebox}

\subsection{El uranio y la energía: la política, los permisos y la geopolítica cambian la curva}

La expansión de la capacidad de cómputo termina dependiendo de la energía, por lo que esta sección sigue el rastro desde los chips y las salas de servidores hasta la generación eléctrica, el combustible y la red eléctrica. El mercado del uranio muestra que el precio de un recurso no depende solo de la oferta de las minas y de la demanda, sino también de la política nuclear, los ciclos de permisos, los inventarios, la geopolítica y la aceptación pública. Aunque una fuente de energía sea técnicamente adecuada para un suministro estable, que pueda llegar a tiempo a los centros de datos sigue dependiendo de la interconexión con la red, el financiamiento, la construcción y la coordinación regulatoria. Al leer la figura hay que distinguir entre el precio del mineral, la electricidad que puede ponerse en operación y la capacidad realmente disponible para un complejo específico; entre los tres hay años de retraso por construcción y aprobación.

\begin{figure}[H]
\centering
\includegraphics[width=0.94\textwidth]{images/00-54-15-uranium-cycle.jpg}
\caption{El ciclo del precio del uranio refleja cómo los recursos, la política, los inventarios y la demanda energética moldean en conjunto el costo de la infraestructura.}
\figsource{Video oficial de Stanford Online 00:54:15, `images/00-54-15-uranium-cycle.jpg`.}
\end{figure}

\begin{knowledgebox}{Lectura de la figura: el precio no es una medida pura de la escasez física}
La cantidad de yacimientos solo determina una parte de la oferta; los permisos de extracción, la capacidad de conversión y enriquecimiento, los contratos de largo plazo, las reservas estratégicas y los planes de construcción de centrales nucleares modifican la oferta disponible. En el debate sobre la energía para la IA también hay que evitar tratar las cifras en GW como electricidad disponible en cualquier momento: la ubicación, el horario, la estabilidad, la capacidad de transmisión y las restricciones de carbono determinan si puede abastecer a un clúster concreto.
\end{knowledgebox}

\subsection{Alternative assets: detrás de una misma forma de onda puede haber mecanismos distintos}

El ponente reúne los ciclos de precios de varias clases de alternative assets para buscar una estructura temporal común que va de la escasez al entusiasmo, a la ampliación de capacidad y a la caída. Esta comparación puede ayudar a los equipos de inversión y construcción a delimitar un rango de escenarios, pero también puede crear con facilidad la ilusión de que «todos los ciclos son iguales». Curvas parecidas pueden deberse a causas distintas, y su duración también depende de la vida útil de los activos, la regulación y los sustitutos.

\begin{figure}[H]
\centering
\includegraphics[width=0.94\textwidth]{images/00-55-45-alternatives-timeline.jpg}
\caption{Comparación de los ciclos de varios alternative assets, para discutir el discurso de la escasez, la afluencia de capital y la respuesta de la oferta.}
\figsource{Video oficial de Stanford Online 00:55:45, `images/00-55-45-alternatives-timeline.jpg`.}
\end{figure}

\begin{warningbox}{Lectura de la figura: una forma parecida no implica la misma causa}
Que un precio suba y luego baje puede deberse a la ampliación de capacidad, al colapso de la demanda, a cambios de política, a la sustitución tecnológica o al desapalancamiento financiero. Si solo se hace una comparación visual, se pasan por alto los plazos de entrega, las diferencias de calidad y los efectos de red de cada activo. Al usar una analogía histórica, primero hay que explicitar el mecanismo y después comprobar si el mercado actual de compute cumple las mismas condiciones.
\end{warningbox}

\begin{table}[H]
\centering
\begin{tabular}{L{0.14\textwidth}L{0.21\textwidth}L{0.23\textwidth}L{0.25\textwidth}}
\toprule
Caso & Mecanismo transferible & Implicación para el compute & Lo que no admite analogía directa \\
\midrule
Acero & La innovación de proceso reduce el costo unitario & La eficiencia del software y de los chips cambia la accesibilidad & La demanda digital crece más rápido \\
Fibra óptica & Tras la sobreconstrucción, los activos depreciados se reutilizan & La capacidad construida por adelantado puede sostener aplicaciones posteriores & Las GPU se deprecian y cambian de generación más rápido \\
DRAM & Desfase entre ciclos largos de ampliación y ciclos cortos de demanda & Los chips, el encapsulado y los inventarios fluctúan & Productos más heterogéneos y cadena de suministro más compleja \\
Uranio & La política, los permisos y la geopolítica moldean la oferta & La electricidad no es un recurso abstracto e ilimitado & La interfaz de entrega de la energía y la del cómputo son distintas \\
Gilded Age & La infraestructura de red trae concentración del mercado & El acceso y el interés público requieren atención institucional & La estructura regulatoria global actual es distinta \\
\bottomrule
\end{tabular}
\caption{El uso correcto de la analogía histórica: transferir mecanismos sin perder de vista las diferencias físicas e institucionales.}
\end{table}

\subsection{Resumen de la sección}

Los ciclos históricos dejan cinco recordatorios: la innovación de proceso reduce el costo unitario, el rebote de la demanda amplía el consumo total, el retraso en la construcción genera ciclos de precios, el valor de la sobreconstrucción depende de que los activos puedan reutilizarse, y la concentración del mercado impulsa los estándares y la regulación pública. La siguiente sección aplica estos recordatorios a un juicio concreto: el compute actual todavía no puede tratarse como un bien sustituible, como la electricidad.
